% GENERATED FILE. Edit canonical lessons, then run npm run generate:books.
% canonical-source-hash: fnv1a64:a4f9aa9126c371c3
% canonical-lessons: TA-C269-ciranta, TA-C269-vittiyacamana, TA-C269-calippana, TA-C269-eliya, TA-C269-vacatiyana

\chapter{Best, Different, Boring, Simple, Comfortable}
\label{ch:ta-best-different-boring-simple-comfortable}
\hlchaptermodality{269}

\begin{chapteropening}
\textbf{By the end of this chapter:} \emph{I can say best, different, boring, simple and comfortable.}

Complete the last lesson of chapter 269: I can say best, different, boring, simple and comfortable.
\end{chapteropening}

\section[\texorpdfstring{\ta{சிறந்த} (ciṟanta)}{சிறந்த (ciṟanta)}]{\ta{சிறந்த} (ciṟanta) — best, finest}

\label{lesson:TA-C269-ciranta}



\begin{quote}
Before the new one: say the Tamil for gentle, then the Tamil for honest.
\end{quote}

\subsection*{You'll want to know: \ta{சிறந்த}}
\textbf{\ta{சிறந்த}} — \emph{ciṟanta} — \textquotedblleft{}best, finest\textquotedblright{}.

\begin{cousinweb}
\textbf{\ta{சிறந்த}} (\emph{ciṟanta}) is native Dravidian. DEDR 2589 lists Malayalam \emph{ciṟakka}, \textquotedblleft{}to be glorious\textquotedblright{}. It is a participle of \emph{ciṟa}, \textquotedblleft{}to be eminent, to excel\textquotedblright{}.
\end{cousinweb}

\begin{grammarlens}[title={What you've built}]
One more word for this chapter.
\end{grammarlens}

\subsection*{Your turn}
\begin{itemize}
\raggedright
  \item \emph{Say it:} \emph{ciṟanta}
  \item \emph{Say it:} \emph{ciṟanta}, once more
  \item \emph{Say it:} say \emph{ciṟanta}
  \item \emph{Recall:} say the Tamil for gentle, then the Tamil for honest, then say \emph{ciṟanta} again
\end{itemize}

\begin{tcolorbox}[breakable,colback=teal!4,colframe=teal!35!black,title={Before you move on}]
What does \ta{சிறந்த} mean? (\textquotedblleft{}Best, finest\textquotedblright{}.) Say it once more.
\end{tcolorbox}
\section[\texorpdfstring{\ta{வித்தியாசமான} (vittiyācamāṉa)}{வித்தியாசமான (vittiyācamāṉa)}]{\ta{வித்தியாசமான} (vittiyācamāṉa) — different, unusual}

\label{lesson:TA-C269-vittiyacamana}



\begin{quote}
Before the new one: say the Tamil for honest, then the Tamil for best.
\end{quote}

\subsection*{You'll want to know: \ta{வித்தியாசமான}}
\textbf{\ta{வித்தியாசமான}} — \emph{vittiyācamāṉa} — \textquotedblleft{}different, unusual\textquotedblright{}.

\begin{cousinweb}
\textbf{\ta{வித்தியாசமான}} (\emph{vittiyācamāṉa}) is \emph{vittiyācam}, \textquotedblleft{}difference\textquotedblright{}, a well-established loan from Sanskrit \emph{vyatyāsa}, \textquotedblleft{}inversion, contrast\textquotedblright{}, with \emph{-āṉa}.
\end{cousinweb}

\begin{grammarlens}[title={What you've built}]
One more word for this chapter.
\end{grammarlens}

\subsection*{Your turn}
\begin{itemize}
\raggedright
  \item \emph{Say it:} \emph{vittiyācamāṉa}
  \item \emph{Say it:} \emph{vittiyācamāṉa}, once more
  \item \emph{Say it:} say \emph{vittiyācamāṉa}
  \item \emph{Recall:} say the Tamil for honest, then the Tamil for best, then say \emph{vittiyācamāṉa} again
\end{itemize}

\begin{tcolorbox}[breakable,colback=teal!4,colframe=teal!35!black,title={Before you move on}]
What does \ta{வித்தியாசமான} mean? (\textquotedblleft{}Different, unusual\textquotedblright{}.) Say it once more.
\end{tcolorbox}
\section[\texorpdfstring{\ta{சலிப்பான} (calippāṉa)}{சலிப்பான (calippāṉa)}]{\ta{சலிப்பான} (calippāṉa) — boring}

\label{lesson:TA-C269-calippana}



\begin{quote}
Before the new one: say the Tamil for best, then the Tamil for different.
\end{quote}

\subsection*{You'll want to know: \ta{சலிப்பான}}
\textbf{\ta{சலிப்பான}} — \emph{calippāṉa} — \textquotedblleft{}boring\textquotedblright{}.

\begin{cousinweb}
\textbf{\ta{சலிப்பான}} (\emph{calippāṉa}) is native: \emph{calippu}, \textquotedblleft{}weariness\textquotedblright{} (DEDR 2371, which lists Malayalam calippu), with \emph{-āṉa}.
\end{cousinweb}

\begin{grammarlens}[title={What you've built}]
One more word for this chapter.
\end{grammarlens}

\subsection*{Your turn}
\begin{itemize}
\raggedright
  \item \emph{Say it:} \emph{calippāṉa}
  \item \emph{Say it:} \emph{calippāṉa}, once more
  \item \emph{Say it:} say \emph{calippāṉa}
  \item \emph{Recall:} say the Tamil for best, then the Tamil for different, then say \emph{calippāṉa} again
\end{itemize}

\begin{tcolorbox}[breakable,colback=teal!4,colframe=teal!35!black,title={Before you move on}]
What does \ta{சலிப்பான} mean? (\textquotedblleft{}Boring\textquotedblright{}.) Say it once more.
\end{tcolorbox}
\section[\texorpdfstring{\ta{எளிய} (eḷiya)}{எளிய (eḷiya)}]{\ta{எளிய} (eḷiya) — simple, easy}

\label{lesson:TA-C269-eliya}



\begin{quote}
Before the new one: say the Tamil for different, then the Tamil for boring.
\end{quote}

\subsection*{You'll want to know: \ta{எளிய}}
\textbf{\ta{எளிய}} — \emph{eḷiya} — \textquotedblleft{}simple, easy\textquotedblright{}.

\begin{cousinweb}
\textbf{\ta{எளிய}} (\emph{eḷiya}) is native Dravidian. DEDR 855 lists Malayalam \emph{eḷutu}, \textquotedblleft{}easy\textquotedblright{}. It comes from \emph{eḷi}, \textquotedblleft{}easy\textquotedblright{}.
\end{cousinweb}

\begin{grammarlens}[title={What you've built}]
One more word for this chapter.
\end{grammarlens}

\subsection*{Your turn}
\begin{itemize}
\raggedright
  \item \emph{Say it:} \emph{eḷiya}
  \item \emph{Say it:} \emph{eḷiya}, once more
  \item \emph{Say it:} say \emph{eḷiya}
  \item \emph{Recall:} say the Tamil for different, then the Tamil for boring, then say \emph{eḷiya} again
\end{itemize}

\begin{tcolorbox}[breakable,colback=teal!4,colframe=teal!35!black,title={Before you move on}]
What does \ta{எளிய} mean? (\textquotedblleft{}Simple, easy\textquotedblright{}.) Say it once more.
\end{tcolorbox}
\section[\texorpdfstring{\ta{வசதியான} (vacatiyāṉa)}{வசதியான (vacatiyāṉa)}]{\ta{வசதியான} (vacatiyāṉa) — comfortable, convenient}

\label{lesson:TA-C269-vacatiyana}



\begin{quote}
Before the new one: say the Tamil for boring, then the Tamil for simple.
\end{quote}

\subsection*{You'll want to know: \ta{வசதியான}}
\textbf{\ta{வசதியான}} — \emph{vacatiyāṉa} — \textquotedblleft{}comfortable, convenient\textquotedblright{}.

\begin{cousinweb}
\textbf{\ta{வசதியான}} (\emph{vacatiyāṉa}) is \emph{vacati}, \textquotedblleft{}comfort, means\textquotedblright{}, usually traced to Sanskrit \emph{vasati}, \textquotedblleft{}dwelling, abode\textquotedblright{}, with \emph{-āṉa}.
\end{cousinweb}

\begin{grammarlens}[title={What you've built}]
One more word for this chapter.
\end{grammarlens}

\subsection*{Your turn}
\begin{itemize}
\raggedright
  \item \emph{Say it:} \emph{vacatiyāṉa}
  \item \emph{Say it:} \emph{vacatiyāṉa}, once more
  \item \emph{Say it:} say \emph{vacatiyāṉa}
  \item \emph{Recall:} say the Tamil for boring, then the Tamil for simple, then say \emph{vacatiyāṉa} again
\end{itemize}

\begin{tcolorbox}[breakable,colback=teal!4,colframe=teal!35!black,title={Before you move on}]
What does \ta{வசதியான} mean? (\textquotedblleft{}Comfortable, convenient\textquotedblright{}.) Say it once more.
\end{tcolorbox}
